\subsection{The exact first nonzero integral weight of the norm obstruction}

\begin{statement}{A.195}{Proved}
{Weight two detects the explicit norm phantom}
Retain the actual obstruction of A.181--A.183.  Thus
\(k=\Q(\alpha)\), \(\alpha^3=3\), \(d=649\),
\(A=d^8\alpha\), \(U=\alpha^2-2\), and \(B=U^3\).
Let \(S\) be the places above \(2,3,11,59,\infty\),
\(\Gamma=G_{k,S}\), and let \(D_s=D_A/T^sD_A\)
with exactly the action and coefficient maps of A.183.
Write \(\mathfrak o_{S,s}\) for the supported integral
obstruction and \(\mathfrak o_s\) for its inflation to
the absolute Galois group.  Then
\[
             \mathfrak o_1=0,\qquad \mathfrak o_2\ne0.
\tag{A.195.1}
\]
In particular the least detecting integral weight is exactly two.
The stronger finite supported assertion is
\[
 \mathfrak o_{S,2}\bmod9\ne0
       \quad\hbox{in }H^2(\Gamma,D_2/9).
\tag{A.195.2}
\]
The supported image modulo three is zero: its actual
coefficient map is the connecting image of
\(\kappa_3(B)=\kappa_3(U^3)=0\).  Thus nine is the
first detecting supported coefficient modulus at weight two.
Its image in \(H^2(G_k,D_2/9)\) is zero, as are all
absolute finite-coefficient images of \(\mathfrak o_2\).
Consequently the coefficient-power inverse-limit comparison of
A.183 places this particular nonzero class in
\[
 \ker\left(H^2_{\rm cts}(G_k,D_2)
                  \longrightarrow\varprojlim_c
                         H^2(G_k,D_2/3^c)\right)
 =\operatorname{im}\left(\varprojlim\nolimits_c^1
                      H^1(G_k,D_2/3^c)\right).
\tag{A.195.3}
\]
All supported and absolute maps in these assertions are the
actual coefficient and inflation maps.
\end{statement}

\begin{proof}
\emph{The supported permutation module and its complete norm kernel.}
Put \(X=1+T\), and set
\[
 P_9=(\Z/9)[T]/T^3,
 \qquad P_3=P_9/3,
 \qquad M=TP_3.
\tag{A.195.4}
\]
On \(P_9\) use the actual \(V_A\)-action
\(g f=\chi(g)X^{a(g)}f(X^{\chi(g)}-1)\), where
\(a=\kappa_9(A)\), \(\chi=\chi_3\bmod9\).
Reduction identifies \(P_3\) with the permutation module
on the three roots of \(A\), with \(\mu_3\) coefficients:
indeed \(X^3-1=T^3\) in characteristic three.
Its root field is the actual \(L=K_2=k(t)\),
\(t^3=d^2\alpha\), since \(A^{1/3}=d^2t\).
The augmentation sequence is therefore
\[
 0\longrightarrow M\longrightarrow
       \operatorname{Ind}_{G_{L,S_L}}^{\Gamma}\mu_3
       \xrightarrow{\ \varepsilon\ }\mu_3\longrightarrow0.
\tag{A.195.5}
\]
The map on first cohomology is exactly the field norm on
the supported Kummer classes, by the Shapiro coset-sum formula.

A.165 gives the actual S-unit bases of dimensions eight and
twelve over \(k\) and \(L\), respectively.  Their localized
class groups have zero three-part by A.177--A.179, so these
are the complete supported Kummer groups.  Here is the
exact norm rank, retaining the displayed generators of A.165.
The four global unit norms have exponents divisible by three.
The eight remaining norms, modulo cubes, lie in the span of
\[
 \begin{gathered}
 [\beta_1],\quad
 [\alpha]-2(k_6+k_7+k_8)[U],\quad
 [\delta_{11}]+k_6[U],\\
 [g_{11}g_{59}]+k_7[U],\quad
 [\delta_{59}]+k_8[U],
 \qquad\delta_q=q/g_q.
 \end{gathered}
\tag{A.195.6}
\]
Each of these five vectors occurs as a norm.  The other
three columns are twice the first vector or zero.  The
five displayed vectors are independent, as seen from their
\(\beta_1,\alpha,\delta_{11},g_{11},\delta_{59}\)
coordinates in the actual eight-element base S-unit basis.
Thus the supported norm has rank five.

Restriction from \(k\) to \(L\) has the one-dimensional
kernel \(\F_3[A]\).  For precision, after the quadratic
extension \(k(\mu_3)/k\) this is ordinary Kummer theory
for the cyclic cubic field obtained by adjoining \(A^{1/3}\).
Restriction through a quadratic extension is injective on
classes modulo cubes, because its norm composite is squaring.
This proves the asserted kernel already over \(k\).
Roots of supported units remain supported units, so this
is the same statement on the displayed S-unit groups.
Restriction consequently has rank seven; its norm composite
is multiplication by three and is zero.  The norm kernel
also has dimension \(12-5=7\).  Therefore
\[
 \ker\bigl(N:H^1(G_{L,S_L},\mu_3)
                    \to H^1(\Gamma,\mu_3)\bigr)
       =\operatorname{res}H^1(\Gamma,\mu_3).
\tag{A.195.7}
\]

Neither \(k\) nor \(L\) contains \(\mu_3\).  The long
exact sequence of (A.195.5) therefore injects \(H^1(\Gamma,M)\)
into the middle first cohomology group, with precisely that
norm kernel as image.  In the basis \(1,X,X^2\), restriction
is the diagonal map, whose vector is
\(1+X+X^2=T^2\) modulo three.  Hence the actual submodule
map
\[
 \mu_3\longrightarrow M,\qquad c\longmapsto T^2c
 \quad\text{induces a surjection on }H^1(\Gamma,-).
\tag{A.195.8}
\]
This is the complete supported source, not just a collection
of local test vectors.

\emph{The coefficient Bockstein and its bottom projection.}
Multiplication by \(T\) identifies \(D_2/3\) with \(M\).
Under this identification the coefficient sequence is
\[
 0\longrightarrow M\xrightarrow{\ j\ }D_2/9
       \xrightarrow{\ r\ }M\longrightarrow0,
 \qquad j(m)=3m/T,\quad r(h)=T(h\bmod3).
\tag{A.195.9}
\]
These maps are equivariant by the defining identity between
the \(D_A\)-action and multiplication by \(T\) in \(V_A\).
Let \(\partial\) be its connecting map from \(H^1(\Gamma,M)\)
to \(H^2(\Gamma,M)\).  The last line \(\F_3T^2\subset M\)
is \(\mu_3\), while
\[
                M/\F_3T^2\simeq\F_3(2).
\tag{A.195.10}
\]
In \(D_2/9\) the corresponding actual stable line is
\(T(\Z/9)(3)\).  Thus naturality of the two coefficient
sequences, together with (A.195.8), proves
\[
 \operatorname{im}\partial
       \subset\operatorname{im}
          \bigl(H^2(\Gamma,\mu_3)\to H^2(\Gamma,M)\bigr).
\tag{A.195.11}
\]
In particular every supported coefficient Bockstein has zero
image in \(H^2(\Gamma,\F_3(2))\).  This also follows directly:
lift a last-line class by \(T\kappa_9(v)\); its coboundary
divided by three stays in the last line.

\emph{The literal augmentation defect.}
Write \(\psi_g=X^{a(g)}\theta_{\chi(g)}\) for the
unweighted root action and put
\[
 w=1+X+X^2=3+3T+T^2\in P_9,
 \qquad \vartheta(g)=\frac{\chi(g)^2-1}{3}\bmod3.
\tag{A.195.12}
\]
The difference \(\psi_gw-w\) is divisible by three.
Its quotient defines an \(I_3=M(-1)\)-valued cocycle.
The exact coefficient identity is
\[
 \frac{\psi_gw-w}{3}
       =(\psi_gX-X)+\vartheta(g)T^2
                      \quad\text{in }I_3.
\tag{A.195.13}
\]
Indeed the two coefficients on the left are
\(a+\chi-1\) and
\(\binom a2+a\chi+\binom\chi2+(\chi^2-1)/3\).
The equality follows from
\(\binom{a+\chi}{2}=\binom a2+a\chi+\binom\chi2\).
Changing integral lifts of \(a,\chi\) does not change
these expressions modulo three.

Take the actual finite Kummer cocycle \(u_9=\kappa_9(U)\).
The continuous cochain \(z(g)=wu_9(g)\) has augmentation
\(3u_9=\kappa_9(B)\).  Its coboundary, followed by division
by \(T\), gives the supported obstruction in \(D_2/9\).
Equation (A.195.13) proves precisely
\[
 \mathfrak o_{S,2}\bmod9=j_*(\gamma),\qquad
 \gamma=\delta_N\kappa_3(U)
       +i_*\bigl(\vartheta\cup\kappa_3(U)\bigr),
\tag{A.195.14}
\]
where \(\delta_N\) is the connecting map of (A.195.5)
and \(i:\mu_3\to M\) is multiplication by \(T^2\).
For the sign, the left-cocycle convention gives
\(dz(g,h)=(g w-\chi(g)w)u_9(h)\), which is three
times the displayed cup cocycle.  The term
\(\psi_gX-X\) is the connecting cocycle obtained by
lifting augmentation one with \(X\), so it gives
the positive \(\delta_N\) in (A.195.14).

Projecting (A.195.14)'s class before \(j\) through
(A.195.10) gives
\[
          \overline\gamma
                 =\kappa_3(A)\cup\kappa_3(U).
\tag{A.195.15}
\]
The possible term \(\chi-1\) from using \(X\) rather
than one as augmentation lift is the coboundary of one
in \(\mu_3\), so has zero cup class.  The extra
\(\vartheta\)-term is in the last line and has zero
bottom projection; no separate vanishing of that term
is required.

The class in (A.195.15) is nonzero.  At the quadratic
prime of \(k\) above eleven (also at the corresponding
prime above fifty-nine), the residue field is \(\F_{q^2}\),
\(v(A)=8\), and the literal cubic residue character
of \(U\) is two in the normalization of A.159.
The local field contains \(\mu_3\); its cubic extension
by \(A^{1/3}\) is totally ramified.  A unit norm has
cube residue, whereas \(U\) does not.  The Kummer
cup/cyclic-algebra map therefore gives a nonzero local
class.  These residue computations are checked again
in the certificate below.

By (A.195.11), \(\gamma\) cannot belong to
\(\operatorname{im}\partial\).  Exactness of (A.195.9)
then gives \(j_*\gamma\ne0\), proving (A.195.2).

\emph{Integral inflation and the absolute finite images.}
The nonzero finite supported image makes
\(\mathfrak o_{S,2}\ne0\) integrally.  The actual
torsion-free inflation map
\[
 H^2(\Gamma,D_2)\longrightarrow H^2_{\rm cts}(G_k,D_2)
\tag{A.195.16}
\]
is injective by the continuous-cochain proof of A.183:
a cobounding cochain on \(G_k\) restricts to an equivariant
homomorphism on the kernel; each outside-S inertia image
is zero because the two diagonal Frobenius weights on
\(D_2(-1)\) are \(Q\) and \(Q^2\), neither one.
The kernel is normally generated by these inertias, so
the cochain descends.  This proves \(\mathfrak o_2\ne0\).
A.183 already proves \(\mathfrak o_1=0\).

At finite coefficients the inflation of (A.195.2) is
zero, by the actual A.181 coefficient map through
\(D_A/(9,T^2)\).  Thus no finite absolute nonvanishing
has been used.  The integral coefficient inverse-limit
sequence proved in A.183 gives (A.195.3) for this
explicitly detected second-weight class.

The exact finite module identities, all twenty-seven norm
matrices, and both quadratic-prime residue values are retained
in \path{code/two_inert_least_weight_two_certificate.py},
with its executed output and data.  The complete S-unit
restriction test is additionally retained in
\path{code/two_inert_base_restriction_unit_classes.py}.
\end{proof}
